% ===================================================================
%  TABLEAU N° 11 — La Ville et ses monuments. — L'Incendie.
%  Feuillets 67 a 70 du fac-simile = folios 65 a 68.
%  Correspondance : folio imprime = numero de feuillet - 2.
%  Le feuillet 67 porte anke la ligne d'apparat « TROISIEME SERIE »
%  avant le titre du tableau.
%
%  Ca tabelo esas en du parti : « La Ville et ses monuments » (folios
%  65-67) e « L'Incendie » (folios 67-68). Kom en l'ido korespondanta
%  (texto/io/20-tabelo-11.tex), la duesma parto riprenas sua propra
%  numerado d'alineas de 1 — videz §6 dil consigne pri la clefi %%K.
%  Kontree l'ido, la francais ne portas nula intertitre decorativa
%  interne (« La Preurbo. », « La Maxim Alta Parto dil Urbo. »,
%  « La Transportili. », « Arto ed Instrukto. », « Klamo pri
%  Incendio. », « Brava Prodaji. ») : nur « L'Incendie. » esas komuna.
%
%  ATTENTION — LE FAC-SIMILE FRANCAIS EST UNE PRISE DE VUE, non un
%  passage au scanner : l'echelle change d'une page a l'autre.
% ===================================================================

% --- Feuillet 67 (folio 65, ouverture de tableau) -------------------
\begin{VUpage}[67]{65}
%%K t11-apar-1 apar
\VUsaut{2.0mm}
\VUcentre{13.2pt}{90}{TROISIÈME SÉRIE}
\VUsaut{3.0mm}
\VUfilet{16mm}
\VUsaut{5.6mm}
\VUtitre{15.0pt}{60}{TABLEAU N\textsuperscript{o} 11}
\VUsaut{1.4mm}
\VUtitre{11.4pt}{0}{La Ville et ses monuments. --- L'Incendie.}
\VUsaut{2.2mm}

%%K t11-c1-01-1 p
\VUblancAlinea
1. --- J'habite une grande ville située à 100 kilo\cc
mètres de l'Océan et qui est pourtant un \VUgras{port}\textsuperscript{(1)} de\nl
premier ordre. Le fleuve qui l'arrose est large de\nl
400 mètres et il forme une très belle rade.

%%K t11-c1-02-1 p
\VUblancAlinea
2. --- Toute la partie de la ville qui se trouve sur les\nl
\VUgras{quais}\textsuperscript{(2)} a un aspect absolument maritime et indus\cc
triel et forme un \VUgras{faubourg}\textsuperscript{(3)} qui n'est habité que par\nl
des marins et par des ouvriers. Ceux-ci travaillent aux\nl
\VUgras{fabriques}\textsuperscript{(4)} et à l'\VUgras{usine à gaz}\textsuperscript{(5)} dont la haute\nl
\VUgras{cheminée}\textsuperscript{(6)} et le \VUgras{gazomètre} se voient de très loin.

%%K t11-c1-03-1 p
\VUblancAlinea
3. --- Au moyen âge, la ville était fortifiée, et l'on\nl
rencontre encore quelques traces de ses remparts, en\nl
particulier la \VUgras{porte}\textsuperscript{(8)} du vieux \VUgras{château-fort}\textsuperscript{(7)}\nl
flanqué de ses deux \VUgras{tours}\textsuperscript{(9)}, qui servent aujourd'hui\nl
de \VUgras{prison.}

%%K t11-c1-04-1 p
\VUblancAlinea
4. --- Dans tout ce \VUgras{quartier,} dont l'aspect est très\nl
ancien, un seul bâtiment est relativement moderne :\nl
c'est la \VUgras{douane}\textsuperscript{(10)}. Elle est placée entre le quai et\nl
la petite \VUgras{rue}\textsuperscript{(11)} qui conduit au vieil \VUgras{hôtel de ville}\textsuperscript{(12)}.\nl
On reconnaît facilement ce dernier monument grâce\nl
au gigantesque \VUgras{beffroi}\textsuperscript{(13)} qui domine toute la vieille\nl
cité.

%%K t11-c1-05-1 p
\VUblancAlinea
5. --- De l'autre côté de la rue s'élève un bâtiment\nl
construit dans le même style que la douane : c'est la\nl
\VUgras{bourse}\textsuperscript{(14)}. L'une de ses façades donne sur une\nl
petite place où se trouve la \VUgras{préfecture}\textsuperscript{(15)}. Notre\nl
ville, en effet, sans être une capitale, est cependant un\nl
important chef-lieu de département.
\end{VUpage}

% --- Feuillet 68 (folio 66) ------------------------------------------
\begin{VUpage}[68]{66}
%%K t11-c1-06-1 p
\VUblancAlinea
6. --- La garnison, assez considérable, est logée\nl
dans de vastes \VUgras{casernes}\textsuperscript{(16)} très saines; elles s'éten\cc
dent jusqu'aux grilles du \VUgras{parc municipal}\textsuperscript{(17)}. Le\nl
\VUgras{boulevard}\textsuperscript{(19)} qui conduit à ce parc passe derrière la\nl
\VUgras{caisse d'épargne}\textsuperscript{(18)} et aboutit à la place de la\nl
\VUgras{cathédrale}\textsuperscript{(20)}.

%%K t11-c1-07-1 p
\VUblancAlinea
7. --- Tous ces monuments s'étalent en amphithéâtre\nl
sur une petite colline où se trouve également notre\nl
vaste \VUgras{lycée}\textsuperscript{(21)}.

%%K t11-c1-07-2 p
Si l'on redescend par une voie un peu en pente qui\nl
rejoint la Grande-Rue, on arrive au \VUgras{palais de jus}\cc
\VUgras{tice}\textsuperscript{(22)} devant lequel s'étend un agréable \VUgras{jardin}\nl
\VUgras{public}\textsuperscript{(26)}. Ce \VUgras{square} n'est pas très grand, mais il\nl
met au milieu de la ville une oasis de verdure et de\nl
fraîcheur. Aussi est-il très fréquenté par les habitants\nl
de la ville, qui aiment à venir s'asseoir sous la tente\nl
du \VUgras{café}\textsuperscript{(25)} pour écouter la musique militaire qui joue\nl
le jeudi et le dimanche dans le petit \VUgras{kiosque}\textsuperscript{(27)} placé\nl
au centre des pelouses et des massifs.

%%K t11-c1-08-1 p
\VUblancAlinea
8. --- Sur une de ces pelouses s'élève une \VUgras{statue}\textsuperscript{(28)}\nl
de bronze, monument érigé à la mémoire de l'un de\nl
nos compatriotes qui a rendu de grands services à sa\nl
ville natale.

%%K t11-c1-09-1 p
\VUblancAlinea
9. --- Notre ville n'est pas encore éclairée à la lumière\nl
électrique; elle n'a que des \VUgras{becs de gaz}\textsuperscript{(29)}. Mais\nl
les \VUgras{journaux} locaux font une campagne pour que ce\nl
système d'éclairage soit amélioré, comme on a déjà\nl
perfectionné le système de traction des tramways.\nl
Les anciennes voitures, traînées par des chevaux, ont\nl
été remplacées par de pratiques \VUgras{tramways électri}\cc
\VUgras{ques}\textsuperscript{(31)} qui transportent rapidement les voyageurs\nl
d'un bout de la ville à l'autre, pour un prix très\nl
modique.

%%K t11-c1-10-1 p
\VUblancAlinea
10. --- Près du palais de justice se trouve la \VUgras{ban}\cc
\VUgras{que}\textsuperscript{(32)}, où l'on escompte les valeurs commerciales.\nl
Les \VUgras{garçons de banque}\textsuperscript{(33)} vont à domicile faire les\nl
recouvrements.
\end{VUpage}

% --- Feuillet 69 (folio 67) ------------------------------------------
\begin{VUpage}[69]{67}
%%K t11-c1-11-1 p
\VUblancAlinea
11. --- Le bel édifice qui se trouve à côté est le\nl
\VUgras{musée}. Il contient à la fois la \VUgras{bibliothèque}\textsuperscript{(34)} de la\nl
ville, de belles \VUgras{collections d'histoire naturelle}\textsuperscript{(35)}\nl
(muséum) et une jolie \VUgras{galerie de tableaux}\textsuperscript{(36)} qui\nl
constitue une superbe \VUgras{exposition} permanente.

%%K t11-c1-11-2 p
Il est ouvert au public deux fois par semaine, mais\nl
les étrangers peuvent le visiter tous les jours en\nl
donnant un pourboire au \VUgras{gardien}\textsuperscript{(37)}. Les \VUgras{artistes}\nl
\VUgras{peintres}\textsuperscript{(38)} viennent y travailler. On les voit devant\nl
leur chevalet, la \VUgras{palette} à la main, recopiant les\nl
tableaux célèbres.

%%K t11-c1-12-1 p
\VUblancAlinea
12. --- Sur la hauteur, derrière le musée, on aper\cc
çoit le \VUgras{dôme}\textsuperscript{(40)} de l'\VUgras{hôpital}\textsuperscript{(39)} et les bâtiments de\nl
l'\VUgras{école normale}\textsuperscript{(41)} où sont formés les instituteurs\nl
de nos écoles communales. Sur la place du Théâtre se\nl
trouve un \VUgras{bureau de poste}\textsuperscript{(43)} installé dans une\nl
\VUgras{maison particulière}\textsuperscript{(42)}.

%%K t11-tit-1 sub
\VUsaut{5.0mm}
\VUpk{11.4pt}{40}{L'Incendie.}
\VUsaut{3.0mm}

%%K t11-c2-01-1 p
\VUblancAlinea
1. --- Un bruit sinistre se répand dans la ville : le feu\nl
vient d'éclater au théâtre. Au moment où j'arrive sur\nl
la \VUgras{place}\textsuperscript{(96)} la foule est déjà massée sur le \VUgras{trot}\cc
\VUgras{toir}\textsuperscript{(46)} et sur la \VUgras{chaussée}\textsuperscript{(47)}. Les \VUgras{employés des}\nl
\VUgras{postes}\textsuperscript{(44)} et les \VUgras{facteurs}\textsuperscript{(45)} sortent du bureau. Un\nl
\VUgras{wattman}\textsuperscript{(48)} a été obligé d'arrêter son tramway\nl
pour laisser passer une \VUgras{voiture des ambulances}\nl
\VUgras{urbaines}\textsuperscript{(50)} qui arrive avec un \VUgras{chirurgien}\textsuperscript{(52)} et\nl
un \VUgras{infirmier}\textsuperscript{(51)} pour prendre un \VUgras{blessé}\textsuperscript{(53)} que l'on\nl
a apporté sur un \VUgras{brancard}\textsuperscript{(54)} auprès du \VUgras{kiosque}\nl
\VUgras{à journaux}\textsuperscript{(55)}.

%%K t11-c2-02-1 p
\VUblancAlinea
2. --- Une compagnie d'infanterie, qui revenait de\nl
l'exercice, s'est arrêtée pour assurer le service d'or\cc
dre avec les \VUgras{gardes municipaux}\textsuperscript{(61)} à cheval. Les\nl
\VUgras{cavaliers,} dont les chevaux se cabrent, font mieux\nl
reculer les \VUgras{curieux}\textsuperscript{(56)} que les \VUgras{soldats}\textsuperscript{(57)} ou les\nl
\VUgras{gendarmes}\textsuperscript{(63)}. D'ailleurs, ces derniers sont fort\nl
occupés. L'un d'eux indique aux \VUgras{agents de police}\textsuperscript{(64)}
\parplein
\end{VUpage}

% --- Feuillet 70 (folio 68) ------------------------------------------
\begin{VUpage}[70]{68}
%%K t11-c2-02-1 p suite
\VUcontinue
où il faut transporter une nouvelle \VUgras{victime}\textsuperscript{(65)}, un\nl
brave sapeur-pompier tombé d'une échelle.

%%K t11-c2-03-1 p
\VUblancAlinea
3. --- L'\VUgras{officier}\textsuperscript{(66)} précise l'endroit où doit se\nl
placer la \VUgras{pompe à vapeur}\textsuperscript{(67)} qui arrive. En atten\cc
dant cette puissante machine, il a fait mettre en\nl
batterie une \VUgras{pompe}\textsuperscript{(68)} à bras, dont le long \VUgras{tuyau}\textsuperscript{(69)}\nl
déverse sur le brasier des torrents d'eau. Six pompiers\nl
la manœuvrent pendant que leur camarade, qui tient\nl
la \VUgras{lance}\textsuperscript{(70)}, dirige le \VUgras{jet}\textsuperscript{(71)} sur le foyer de l'incendie.

%%K t11-c2-04-1 p
\VUblancAlinea
4. --- De l'autre côté du théâtre, on a dressé la\nl
grande \VUgras{échelle}\textsuperscript{(72)}. Elle permet aux sapeurs de s'ap\cc
procher des flammes qui jaillissent au milieu de tour\cc
billons de \VUgras{fumée}\textsuperscript{(75)} noire.

%%K t11-c2-05-1 p
\VUblancAlinea
5. --- Le \VUgras{feu}\textsuperscript{(74)} a pris naissance dans le foyer du\nl
\VUgras{théâtre}\textsuperscript{(73)}, pendant une répétition de l'\VUgras{opéra}\nl
dont la première \VUgras{représentation} devait avoir lieu\nl
demain. Il n'y avait heureusement que peu de \VUgras{spec}\cc
\VUgras{tateurs}\textsuperscript{(76)} dans la salle. Les malheureux se sont\nl
réfugiés sur le \VUgras{balcon}\textsuperscript{(77)}, où de courageux \VUgras{pom}\cc
\VUgras{piers}\textsuperscript{(86)} viennent à leur secours avec une échelle et\nl
des cordages.

%%K t11-c2-06-1 p
\VUblancAlinea
6. --- Sous le \VUgras{péristyle}\textsuperscript{(78)}, où l'\VUgras{affiche}\textsuperscript{(79)} annonce\nl
la représentation, on voit les \VUgras{musiciens}\textsuperscript{(81)} de\nl
l'\VUgras{orchestre}\textsuperscript{(80)} s'enfuir en emportant leurs instru\cc
ments : \VUgras{violon}\textsuperscript{(81)}, \VUgras{flûte}\textsuperscript{(82)}, \VUgras{violoncelle}\textsuperscript{(83)}, \VUgras{trom}\cc
\VUgras{bone}\textsuperscript{(84)}, \VUgras{tambour}\textsuperscript{(85)}, etc. On essaie de sauver une\nl
partie du \VUgras{mobilier}\textsuperscript{(87)}.

%%K t11-c2-07-1 p
\VUblancAlinea
7. --- Les \VUgras{acteurs}\textsuperscript{(89)} et les \VUgras{actrices}\textsuperscript{(88)}, \VUgras{chan}\cc
\VUgras{teurs} et \VUgras{cantatrices,} ont dû fuir avec leur \VUgras{cos}\cc
\VUgras{tume}\textsuperscript{(90)} de théâtre. Le ténor explique à un \VUgras{jour}\cc
\VUgras{naliste}\textsuperscript{(92)} comment les choses se sont passées. Le\nl
\VUgras{reporter} est accouru dès le premier moment avec\nl
les autorités : le \VUgras{préfet}\textsuperscript{(91)}, le \VUgras{maire}\textsuperscript{(93)}, le \VUgras{com}\cc
\VUgras{missaire de police}\textsuperscript{(94)} et un \VUgras{magistrat}\textsuperscript{(95)}, qui\nl
vont ouvrir une enquête pour établir les responsabi\cc
lités.

\VUsaut{6.0mm}
\VUfilet{22mm}
\end{VUpage}
